% Ek D -- Hesap notları (revizyon-v2: özgün Ek A.3 ve A.4 ile H05 gradyan normu oranı tanımı)
\setcounter{table}{0}\setcounter{figure}{0}
\section{Yöntem ayrıntıları ve hesap notları}\label{sec:ek-hesap}

Bu ek ana metnin Yöntem bölümünden çıkarılan ayrıntıları (hipotez kartı ve karar kuralları, kayıp terimlerinin açılımı, eğitim reçetesi, müdahale ayrıntıları) ve hesap notlarını verir; ön-eleme olasılıkları, gradyan normu oranı ve makullük çıpasının kaynakları makalenin Ek~\ref{sec:deney-kaydi}'sındadır.

\subsection{Hipotez kartı ve karar kuralı ayrıntıları}\label{sec:ek-kart}
Kart ana metindeki alanlara ek olarak kabul sonrası ablasyonları da sabitler. Kart dosyası ilk koşumdan önce damgalanmalıdır (bu kuralın uygulamadaki kanıt gücü sınırlıdır; Bölüm~\ref{sec:sinirliliklar}). Birleşik kart ancak bileşenlerin her biri ayrı kabul almışsa yazılabilir. Beklenen aralıklar, izole bir VOC ablasyonu bulunamayan kartlarda dardır ve sıfır merkezlidir.

Ön-elemede yanlış eleme maliyeti yüksek olduğundan eşik tek-tohum oynaklığının altında tutulmuştur; sıfır etkide de geçiş olasılığı ihmal edilemediğinden geçiş ikinci doğrulama gerektirir. Sonradan üç tohuma tamamlanan elenen kartta SÜRDÜR ölçütünün karşılanmaması ``isabet'', karşılanması ``kaçırılmış kazanç'' diye betimlenir; isabet bir özgüllük ölçüsüdür ve bilgisiz bir kapıda da olası olduğundan bilgilendirici sayılmaz (makale Ek~\ref{sec:ek-onele}). SÜRDÜR kuralında $s_\Delta$, $\Delta_{k,s}$'lerin örneklem s.d.'sidir ve tek yönlü sınamanın alternatifi $H_1:\bar\Delta_k>0$'dır. İNCELE durumunda ek tohum yalnız K5 bütçesinden ve hem varyant hem taban model için eklenir. SÜRDÜR kuralı tur 2'de, H01--H03 sonuçları görüldükten sonra donmuştur; tur-2 kuralını geçen kart ($\bar\Delta_k\ge+0.40$) tur-1 kuralını da geçer, tersi doğru değildir.

\subsection{Taban model kayıp terimlerinin açılımı}\label{sec:ek-kayip}
Taban modelin boynu üç ölçekte (adım 8, 16, 32) çıktı üreten yol toplama ağı (path aggregation network, PAN) tipindedir (üçüncü taraf bir derleme: \cite{khanam_2024_yolov11overview}). TAL'da pozitiflerin sınıf hedefi $t_i=a_{ij}\max_{i'}u_{i'j}/\max_{i'}a_{i'j}$'dir. Ana metindeki Denklem~\eqref{eq:box}'te geçen CIoU \cite{zheng_2020_diou} ve dağılım odaklı kayıp (distribution focal loss, DFL) \cite{li_2020_gfl} şöyledir:
\begin{align}
\mathrm{CIoU}&=\mathrm{IoU}-\frac{\rho^2(\hat{\mathbf m},\mathbf m)}{d^2}-\alpha_v v,\label{eq:ciou}\\
\mathcal L_{\mathrm{dfl}}&=\frac1T\sum_{i\in\mathcal P}w_i\,\frac14\sum_{e=1}^{4}\Bigl[-\bigl((r_e+1)-y_e\bigr)\log S_{e,r_e}-(y_e-r_e)\log S_{e,r_e+1}\Bigr].\label{eq:dfl}
\end{align}
Denklem~\eqref{eq:ciou}'de $\hat{\mathbf m}$ ve $\mathbf m$ kutu merkezleri, $\rho$ Öklid uzaklığı, $d$ iki kutuyu kapsayan en küçük kutunun köşegeni, $v=\frac4{\pi^2}\bigl(\arctan\frac{w^{gt}}{h^{gt}}-\arctan\frac{\hat w}{\hat h}\bigr)^{2}$ en-boy terimi ve $\alpha_v=v/\bigl((1-\mathrm{IoU})+v\bigr)$'dir \cite{zheng_2020_diou}. BCE, $\mathrm{BCE}(z,t)=-t\log\sigma(z)-(1-t)\log(1-\sigma(z))$'dir. Denklem~\eqref{eq:dfl}'teki DFL \cite{li_2020_gfl} dört kenar uzaklığının her birini $R=16$ ayrık bölme üzerinde bir softmax dağılımı $S_e$ ile tahmin eder; $y_e$ hedef uzaklık, $r_e=\lfloor y_e\rfloor$'dir.

\subsection{Eğitim reçetesi ve veri artırma varsayılanları}\label{sec:ek-recete}
Dönüşümler sırasıyla şunlardır. Olasılığı 1.0 olan mosaic dört görüntüyü $2s\times2s$ boyutlu bir tuvalde birleştirir ($s=640$) \cite{bochkovskiy_2020_yolov4}. RandomPerspective ölçek çarpanını $a\sim U[1-\sigma_s,\,1+\sigma_s]$ ile seçer ($\sigma_s=0.5$), 0.1 oranında öteleme uygular (dönme, kayma ve perspektif sıfırdır) ve çıktıyı merkezdeki $s\times s$ pencereye kırpar. Ardından HSV kayması (hue-saturation-value; 0.015, 0.7, 0.4) ve olasılığı 0.5 olan yatay çevirme gelir. MixUp ve yerleşik copy-paste kapalıdır; mosaic son 10 epoch'ta (\texttt{close\_mosaic=10}) kapanır. Yerleşik copy-paste maske gerektirdiğinden kutu etiketli VOC'ta etkisizdir \cite{ultralytics_2026_augdocs}; H03 ve H08 kutu-tabanlı yeni bir uygulamaya dayanır.

\subsection{Müdahale ayrıntıları}\label{sec:ek-mudahale}
\paragraph{Ölçek ve kutu kaybı.}
H01'in amacı aynı nesneyi farklı piksel yoğunluğunda göstererek ölçek değişmezliğini zorlamaktır. H04'te RandomPerspective ölçek aralığı $[0.5,1.5]$'ten $[0.75,1.25]$'e iner; genlik büyütülmez. H02'de her pozitif çapa için yardımcı kutu $b^{\mathrm{in}}=(x_c,\,y_c,\,rw,\,rh)$'dir; DFL terimi, TAL ve kazançlar değişmez ve uygulama $r=1$ özdeşlik birim testiyle doğrulanmıştır. H05'te etkin gradyan normu ölçülür (makale Ek~\ref{sec:ek-gradyan}).

\paragraph{Kutu-tabanlı copy-paste.}
Kopyala-yapıştır, küçük nesneleri aynı görüntüde çakışmasız çoğaltarak küçük nesne açığını kapatmayı dener \cite{kisantal_2019_smallaug}. Dönüşüm mosaic'ten sonra, RandomPerspective'ten önce çalışır; adaylar başka kutularla örtüşmeyen ($\mathrm{IoA}\le0.1$) nesnelerdir, $k\sim U\{1,2,3\}$ nesne seçilir ve 3 piksellik alfa harmanlamasıyla yapıştırılır (adımlar: Ek~\ref{sec:ek-cp}). Maskesiz yama arka plan parçası da taşır; yama kenarını öğrenme riski AP$_{75}$'in izlenme nedenlerindendir. Dönüşüm son 10 epoch'ta kapalıdır. H03'te $p=0.5$ ve $\tau_A=48$ piksel, H08'de $p=1.0$ ve $\tau_A=96$ pikseldir. Yapıştırma eşikleri mosaic tuvalinde, AP$_S$/AP$_M$/AP$_L$ sınırları orijinal pikselde COCO tanımıyla ölçülür \cite{lin_2014_coco}. Havuzun genişlemesi yapıştırılanların boyut karışımını da değiştirdiğinden H08'de olasılık ve havuz ayrıştırılmamıştır. Doz $D=N_{\text{yapıştırılan}}/N_{\mathrm{mosaic}}$ [nesne/mosaic] işçi süreçlerinin periyodik yazdığı sayaçlardan yaklaşık ölçülür. H08 kartında hedef doz $D\ge0.80$ (beklenen aralık 0.8--1.2) olarak önceden yazılmıştır; koşum öncesi ölçülen $D$, $[0.6,\,1.5]$ kapısının dışındaysa koşum başlamaz. Hedefin altında kalan doz ikinci bir değişken doğurmamak için ayarlanmaz; gerçekleşen doz raporlanır.

\paragraph{Eğitimsiz çıkarım satırları (H07).}
Referans hat R0, 0.001 güven eşiği, 0.7 NMS IoU eşiği, en çok 300 tespit, 640 piksel tek ölçek ve sınıf-bazlı klasik NMS kullanır. Her satır R0'dan yalnız bir bileşenle ayrılır ve satırlar birleştirilmez. \emph{NMS eşiği (N)} IoU eşiğini $\{0.5,\,0.6,\,0.8\}$'den birine alır. \emph{Soft-NMS (S)} örtüşen kutunun skorunu Gauss kuralıyla sınıf başına bozar ($\sigma\in\{0.3,\,0.5\}$; tanım: Ek~\ref{sec:ek-h07-dogrulama}) \cite{bodla_2017_softnms}. \emph{Test çözünürlüğü (U)} girişi 704, 768 ya da 832 piksele çıkarır. \emph{TTA (T)}: T1 Ultralytics varsayılanıdır \cite{ultralytics_2026_yolov5tta,ultralytics_2026_v84165}; ölçekleri $\{1,\,0.83,\,0.67\}$ olan üç dal çalışır, ortadaki dal yatay çevrilir ve dal çıktıları özgün koordinatta tek NMS'e verilir. T2'de ölçek listesi $\{1,\,1.2,\,1.4\}$'tür. Üç dal giriş pikseli toplamını $1+0.83^2+0.67^2\approx2.14$ katına çıkarır; gecikme buna dal birleştirme ve NMS eklenerek artar. \emph{WBF (W)} T1 ya da T2 dallarını ayrı ayrı NMS'ten geçirip füzyonla birleştirir (IoU eşiği 0.6, atlama eşiği 0.001, ortalama güven) \cite{solovyev_2021_wbf}. Gecikme oranı aynı oturumda ölçülen R0'a göre hesaplanır; $\rho\le1.10$ eşiği, gerçek-zaman iddiasının taşınabileceği en büyük gecikme artışı olarak koşumdan önce belirlenmiştir. KAZANÇ satırları eğitim kartlarının ebeveyn zincirine girmez.


\subsection{İstatistik ayrıntıları}\label{sec:ek-istatistik}
Bölüm~\ref{sec:karar}'teki ``İstatistik'' paragrafını ayrıntılandırır. Eğitim kartlarında tohum eşlemesi nominaldir, çünkü müdahale veri artırmanın rastgele sayı üreteci (RNG) tüketimini de değiştirir ve eşleştirme varyansı yalnız kısmen düşürür. Holm düzeltmesinin \cite{holm_1979_sequential} on iki testlik K5 ailesi H03 ve TTA $\times$ altı metrikten oluşur (H03 mAP50-95 dahil, muhafazakâr sayım); Holm-7 yedi kartın üç tohumlu doğrulama farklarına uygulanır. Görüntü-düzeyi eşli bootstrap \cite{efron_1979_bootstrap} aynı görüntü indekslerini varyanta ve taban modele uygular ($B=1000$, yüzde 95 yüzdelik aralığı; üç tohumun $\Delta$ ortalaması). Bu aralık eğitim tohumu varyansını içermediğinden eğitim müdahaleleri için tohum düzeyinde aralık $\bar\Delta_k\pm t_{0.975,2}\,s_\Delta/\sqrt n$ ($t_{0.975,2}=4.303$) de verilir. Düzeltmesiz $p<0.05$ tablolarda renklendirilmez. Pratik marjlar ($+0.30$, $+0.40$ puan) istatistiksel ölçütten ayrı bir ölçüttür. K5'te taban modelin test mAP50-95 değeri her tohum için makullük denetimindeki değerle karşılaştırılır; uyuşmazlık düzeltilmeden raporlanır.

\subsection{Dış kaynak değeri}\label{sec:ek-dis}
KITTI taban modelinin büyüklük mertebesi için tek dış kaynak değeri, bir topluluk modeli kartıdır (YOLO11n, KITTI; \url{https://huggingface.co/dronefreak/kitti-yolo11n}, erişim 2026-10-22): doğrulama mAP50-95'i için 24.0, car AP50-95'i için 67.3 bildirir. Model 1280 px'te ve 3712/3769 bölmesinde eğitilmiş, \texttt{DontCare} işlemesi belirtilmemiş ve erken durdurma (43. epoch) ile raporlanan sayılar aynı doğrulama kümesine dayanır. Çözünürlük, bölme ve değerlendirme farklı olduğundan bizim değerlerimizle karşılaştırılamaz; hakemli bir kaynak değildir.

\FloatBarrier
